\documentclass[12pt]{article}
\usepackage[UTF8]{ctex}
\usepackage[a4paper,margin=1in]{geometry}
\usepackage{amsmath,amssymb,booktabs,array}
\usepackage{enumitem}
\usepackage{multicol}
\usepackage{physics}

\setlength{\parindent}{0pt}
\setlength{\parskip}{0.6em}

\begin{document}

\begin{center}
{\LARGE AP Physics 1: Simple Harmonic Motion and Energy}\\
{\large Practice Set}\\
\end{center}

\section*{Part I. Multiple-Choice Questions}

\subsection*{Question 1}
A block of mass $m$ is attached to an ideal spring of spring constant $k$ on a frictionless horizontal surface. The block is pulled a distance $A$ from its equilibrium position and released from rest. At what distance from the equilibrium position is the block's kinetic energy equal to its spring potential energy?

\begin{multicols}{2}
\begin{enumerate}[label=\textbf{\Alph*.}]
\item $\dfrac{A}{4}$
\item $\dfrac{A}{2}$
\item $\dfrac{A}{\sqrt{2}}$
\item $A$
\end{enumerate}
\end{multicols}

\hrule

\subsection*{Question 2}
A block attached to a spring oscillates with amplitude $A$ on a frictionless horizontal surface. Its maximum speed is $v$.

The block is then replaced with a new block of mass $4m$, and the amplitude is increased to $2A$. The spring constant remains unchanged.

What is the new maximum speed?

\begin{multicols}{2}
\begin{enumerate}[label=\textbf{\Alph*.}]
\item $\dfrac{v}{2}$
\item $v$
\item $2v$
\item $4v$
\end{enumerate}
\end{multicols}

\hrule

\subsection*{Question 3}
A block of mass $m$ is attached to a spring of spring constant $k$ on a horizontal surface. The surface is frictionless except for a rough strip of length $L$ centered at the equilibrium position.

The block is pulled to the right a distance $A$ from equilibrium, where $A>L/2$, and released from rest. It moves left, passes through the rough strip, and momentarily comes to rest on the other side when the spring is compressed by $\dfrac{A}{2}$.

What is the magnitude of the friction force exerted by the rough strip?

\begin{multicols}{2}
\begin{enumerate}[label=\textbf{\Alph*.}]
\item $\dfrac{3kA^2}{8L}$
\item $\dfrac{kA^2}{4L}$
\item $\dfrac{kA^2}{2L}$
\item $\dfrac{3kA^2}{4L}$
\end{enumerate}
\end{multicols}

\hrule

\subsection*{Question 4}
A student measures the period of a simple pendulum released from a very small angle and finds it to be $T_1$. The same pendulum is then released from a much larger angle, and its period is measured to be $T_2$.

Which statement is correct?

\begin{enumerate}[label=\textbf{\Alph*.}]
\item $T_2=T_1$ because the period of all pendulums is independent of amplitude.
\item $T_2>T_1$ because at large angles the restoring force is no longer proportional to displacement, so the motion is not perfectly simple harmonic.
\item $T_2<T_1$ because the pendulum has more gravitational potential energy at release.
\item $T_2>T_1$ only because the pendulum travels a longer arc length.
\end{enumerate}

\hrule

\subsection*{Question 5}
A horizontal spring-block oscillator and a simple pendulum are both moved from Earth to a location where the local gravitational field strength is much smaller.

How do their periods change?

\begin{center}
\begin{tabular}{>{\bfseries}c c c}
\toprule
Choice & Spring-Block Period & Pendulum Period \\
\midrule
A & Increases & Increases \\
B & Unchanged & Increases \\
C & Unchanged & Decreases \\
D & Decreases & Unchanged \\
\bottomrule
\end{tabular}
\end{center}

\hrule

\subsection*{Question 6}
A simple pendulum oscillating with a small amplitude has a period $T$ and a maximum kinetic energy $K_{\max}$. If the amplitude of the pendulum's swing is doubled while remaining small enough for the motion to be considered simple harmonic, what are the new period and new maximum kinetic energy?

\begin{center}
\begin{tabular}{>{\bfseries}c c c}
\toprule
Choice & New Period & New Maximum Kinetic Energy \\
\midrule
A & $T$ & $2K_{\max}$ \\
B & $T$ & $4K_{\max}$ \\
C & $2T$ & $2K_{\max}$ \\
D & $2T$ & $4K_{\max}$ \\
\bottomrule
\end{tabular}
\end{center}

\newpage
\section*{Part II. Free-Response Questions}

\subsection*{Free-Response Question 1: Energy and Motion of a Simple Pendulum}

A simple pendulum consists of a bob of mass $m$ attached to a light string of length $L$. The pendulum is displaced a small angle and released from rest, so that it undergoes simple harmonic motion.

\begin{enumerate}[label=\textbf{(\alph*)}]
\item Explain, in terms of the net force acting on the bob, why the pendulum's motion is considered simple harmonic only for \textit{small} angles.

\item Provide a claim about whether the acceleration of the bob at its maximum displacement depends on the mass $m$ of the bob. Justify your claim using Newton's laws.

\item The zero of gravitational potential energy is defined at the lowest point of the pendulum's path. On separate energy bar charts, indicate the relative amounts of kinetic energy $K$ and gravitational potential energy $U_g$ for each of the following positions:
\begin{enumerate}[label=\textbf{\roman*.}]
\item at the highest point of the motion
\item as the bob passes through the lowest point
\end{enumerate}

\item A student claims that if the bob's mass is doubled while the pendulum is released from the same initial height, then the speed at the bottom of the swing will double. State whether the student is correct, and justify your answer using energy principles.
\end{enumerate}

\hrule

\subsection*{Free-Response Question 2: Determining an Unknown Spring Constant}

A lab group is given an unmarked spring and asked to determine its spring constant $k$. The available equipment includes a set of known masses, a meterstick, a stopwatch, and a support stand.

Design \textbf{two different experiments} to determine $k$.

\begin{enumerate}[label=\textbf{(\alph*)}]
\item Describe a \textbf{static} method that does not involve oscillation.
\item Describe a \textbf{dynamic} method that uses oscillation.
\item For each method, identify the quantities that must be measured directly.
\item For each method, explain how the measured data can be used to calculate $k$.
\item Describe one specific step you would take to reduce experimental uncertainty.
\end{enumerate}

\end{document}